\documentclass[12pt,a4paper]{article}

% ── Packages ──────────────────────────────────────────────────────────────────
\usepackage[margin=1in]{geometry}
\usepackage{amsmath,amssymb}
\usepackage{booktabs}
\usepackage{longtable}
\usepackage{array}
\usepackage{tabularx}
\usepackage{hyperref}
\usepackage{xcolor}
\usepackage{listings}
\usepackage{enumitem}
\usepackage{graphicx}
\usepackage{fancyhdr}
\usepackage{titlesec}
\usepackage{setspace}
\usepackage{parskip}
\usepackage{microtype}
\usepackage{caption}
\usepackage{multirow}

% ── Hyperlink colours ─────────────────────────────────────────────────────────
\hypersetup{
  colorlinks=true,
  linkcolor=black,
  urlcolor=blue!70!black,
  citecolor=black
}

% ── Code listing style ────────────────────────────────────────────────────────
\lstdefinestyle{codestyle}{
  backgroundcolor=\color{gray!8},
  basicstyle=\ttfamily\footnotesize,
  breaklines=true,
  frame=single,
  framerule=0.4pt,
  rulecolor=\color{gray!40},
  xleftmargin=8pt,
  xrightmargin=8pt,
  aboveskip=6pt,
  belowskip=6pt,
  showstringspaces=false,
  keywordstyle=\color{blue!70!black}\bfseries,
  commentstyle=\color{green!50!black},
  stringstyle=\color{red!60!black},
}
\lstset{style=codestyle}

% ── Section formatting ────────────────────────────────────────────────────────
\titleformat{\section}{\large\bfseries}{}{0em}{\thesection.\quad}
\titleformat{\subsection}{\normalsize\bfseries}{}{0em}{\thesubsection\quad}
\titleformat{\subsubsection}{\normalsize\itshape\bfseries}{}{0em}{\thesubsubsection\quad}

% ── Header / footer ───────────────────────────────────────────────────────────
\pagestyle{fancy}
\fancyhf{}
\fancyhead[L]{\small DS Mentor Pro --- Project Report}
\fancyhead[R]{\small Team 2}
\fancyfoot[C]{\thepage}
\renewcommand{\headrulewidth}{0.4pt}

% ── Spacing ───────────────────────────────────────────────────────────────────
\setlength{\parskip}{6pt}
\setlength{\parindent}{0pt}
\onehalfspacing

% ─────────────────────────────────────────────────────────────────────────────
\begin{document}

% ── Title page ────────────────────────────────────────────────────────────────
\begin{titlepage}
  \centering
  \vspace*{2cm}
  {\LARGE\bfseries DS Mentor Pro\par}
  \vspace{0.4cm}
  {\large Agentic Data Science Tutor with Retrieval-Augmented Generation\par}
  \vspace{0.3cm}
  {\normalsize\textit{Project Report}\par}
  \vspace{0.2cm}
  {\small Live Deployment: \href{https://dsmentor.parvpatel.me}{\texttt{dsmentor.parvpatel.me}}\par}
  \vspace{1.5cm}

  {\large\bfseries Team 2\par}
  \vspace{0.8cm}

  \begin{tabular}{ll}
    \toprule
    \textbf{Name} & \textbf{Student ID} \\
    \midrule
    Parv Patel    & 142301041 \\
    Prisha Italiya & 142301011 \\
    Vaibhav Singh & 142502033 \\
    Sangam        & 112201006 \\
    \bottomrule
  \end{tabular}

  \vfill
  {\small\today\par}
\end{titlepage}

\tableofcontents
\newpage

% ─────────────────────────────────────────────────────────────────────────────
\section{Motivation}

Learning data science is inherently sequential: a practitioner who skips Exploratory Data Analysis before Preprocessing, or who applies feature scaling after the train/test split, will produce flawed results. General-purpose language models like ChatGPT answer data science questions in isolation --- they do not track which pipeline stage the user is at, do not detect when a step is skipped, and are prone to hallucination because their answers are not grounded in verified code examples.

Our group chose to improve \textbf{domain-grounded, workflow-aware tutoring} for the data science pipeline. Specifically we target three weaknesses of general LLMs in this domain:

\begin{enumerate}[leftmargin=*, label=\arabic*.]
  \item \textbf{Hallucination} --- LLMs generate plausible-sounding but incorrect code that does not match the user's actual dataset (wrong column names, wrong API signatures).
  \item \textbf{No pipeline awareness} --- LLMs treat every query independently; they never warn the user that they are attempting Stage~6 (Modeling) without completing Stage~4 (Preprocessing).
  \item \textbf{No multi-turn grounding} --- follow-up queries with pronouns (\textit{``how do I fix it?''}, \textit{``show me this for the training set''}) are not resolved against prior context.
\end{enumerate}

Our system, \textbf{DS Mentor Pro}, addresses all three by combining retrieval-augmented generation, a fine-tuned stage classifier, a pipeline tracker, and a multi-turn conversation manager.

% ─────────────────────────────────────────────────────────────────────────────
\section{Why DS Mentor Pro is NOT a Normal Chatbot}

Most QA systems --- including ChatGPT --- operate as stateless question-answering machines: a query goes in, an answer comes out, and each turn is treated independently. DS Mentor Pro is fundamentally different in architecture, behaviour, and purpose.

\begin{table}[h!]
\centering
\caption{DS Mentor Pro vs.\ a Typical QA Chatbot}
\small
\begin{tabularx}{\textwidth}{>{\bfseries}lXX}
\toprule
Capability & Typical QA Chatbot & DS Mentor Pro \\
\midrule
Answer grounding & Generated from model weights (hallucination-prone) & Retrieved from verified DS notebook examples \\
Pipeline awareness & None --- every query is independent & Tracks which of 7 stages the user has completed; warns on skips \\
Stage classification & None & Automatically identifies which DS step a query belongs to \\
Code adaptation & Verbatim generation, may use wrong column names & Extracts column names/models from user's query and adapts code \\
Anti-pattern detection & None & Detects 7 common DS mistakes (data leakage, blind dropna, fit on test, etc.) \\
Multi-turn context & Basic history window & Resolves pronouns (``it'', ``this'', ``the model'') against session entities \\
Visualization & Text description only & Generates and \textit{executes} matplotlib/seaborn code; returns image \\
Confidence scoring & None & Composite score from retrieval quality + syntax validity + critic agent \\
Agentic planning & None & \texttt{PlannerAgent} decomposes complex queries into ordered subtasks \\
Socratic mode & None & Decides when to ask guiding questions instead of direct answers \\
Quiz \& checkpoint & None & Generates MCQ/code quizzes and grades checkpoint submissions per stage \\
Learning analytics & None & SQLite telemetry tracks quiz scores, checkpoints, and per-stage progress \\
Latency & 850\,ms (API roundtrip) & 2.4\,ms retrieval; full response under 50\,ms \\
\bottomrule
\end{tabularx}
\end{table}

The system is deployed live at \href{https://dsmentor.parvpatel.me}{\textbf{dsmentor.parvpatel.me}}, serving real users through a FastAPI backend and a Streamlit interface, with session persistence, report export, and file upload for dataset profiling.

% ─────────────────────────────────────────────────────────────────────────────
\section{Approach}

\subsection{System Architecture}

The system is composed of six interacting modules:

\begin{lstlisting}[language={}]
User Query
    |
    v
+------------------------+
|  Stage Classifier      |  --> stage in {1,...,7}
|  (TF-IDF + SVM)        |
+----------+-------------+
           | stage label
    +------v------------------+
    |  Pipeline Tracker       |  --> skip warnings
    +------+------------------+
           | enriched query
    +------v------------------+
    |  Hybrid Retriever       |  --> top-K docs
    |  (BM25 + TF-IDF + RRF) |
    +------+------------------+
           | retrieved context
    +------v------------------+
    |  Code Generator         |  --> Python snippet
    |  (CodeT5-small + templ.)|
    +------+------------------+
           |
    +------v------------------+
    |  Conversation Manager   |  --> pronoun resolution
    |  + Confidence Scorer    |  --> grounded response
    +-------------------------+
\end{lstlisting}

\subsection{Model Selection Rationale}

Every model in the pipeline was chosen deliberately over more powerful but impractical alternatives.

\subsubsection{Stage Classifier --- TF-IDF + LinearSVC}

\begin{table}[h!]
\centering
\caption{Stage Classifier Options Compared}
\begin{tabular}{lccc}
\toprule
Option & Accuracy & Training Time & Deployment Size \\
\midrule
TF-IDF + LinearSVC & 73.1\% & $\sim$3 seconds & $\sim$2\,MB \\
DistilBERT fine-tune & $\sim$76\% (est.) & $\sim$4 hours (GPU) & $\sim$265\,MB \\
GPT-4 zero-shot & $\sim$80\% (est.) & N/A & API-only \\
\bottomrule
\end{tabular}
\end{table}

\textbf{Why TF-IDF + SVM:} The vocabulary of DS queries is highly domain-specific and keyword-driven (\texttt{fillna}, \texttt{train\_test\_split}, \texttt{fit\_transform}). TF-IDF captures this sparse, discriminative vocabulary exactly. LinearSVC is a max-margin classifier well-suited to high-dimensional sparse features and trains in seconds, making it reproducible without a GPU. The 3\% accuracy gap versus DistilBERT does not justify 80$\times$ larger deployment size for a tutoring tool that must be self-hosted.

\subsubsection{Retriever --- BM25 + TF-IDF + RRF}

We chose a \textbf{sparse hybrid} over dense retrieval (FAISS + sentence-transformers) for three reasons:
\begin{enumerate}[leftmargin=*]
  \item \textbf{No GPU required at inference} --- sparse retrieval runs entirely on CPU.
  \item \textbf{Exact keyword match matters} --- queries like ``how do I use \texttt{StandardScaler}'' benefit from exact BM25 term matching that dense embeddings can miss.
  \item \textbf{RRF outperforms either signal alone} --- our grid-search confirms this at no extra cost.
\end{enumerate}

\subsubsection{Code Generator --- CodeT5-small}

\begin{table}[h!]
\centering
\caption{Code Generator Options Compared}
\begin{tabular}{lcll}
\toprule
Option & CodeBLEU & Training Feasibility & Notes \\
\midrule
Base CodeT5-small (no fine-tune) & 31\% & N/A & Lowest baseline \\
\textbf{Fine-tuned CodeT5-small} & \textbf{69.2\%} & Yes (1--3\,hrs GPU) & \textbf{Our choice} \\
CodeT5-plus (220M) & $\sim$75\% (est.) & Needs A100 & Too large for our setup \\
ChatGPT-3.5 (zero-shot) & 52\% & No training & Paid API; hallucination risk \\
\bottomrule
\end{tabular}
\end{table}

\subsection{Stage Classifier (Mathematical Formulation)}

Given a query $q$, let $\mathbf{x} = \text{TF-IDF}(q) \in \mathbb{R}^V$. The classifier predicts:
\[
  \hat{s} = \arg\max_{s \in \{1,\ldots,7\}} \bigl(\mathbf{w}_s \cdot \mathbf{x} + b_s\bigr)
\]
where $\mathbf{w}_s, b_s$ are the LinearSVC weights and bias for stage $s$, learned by minimizing hinge loss:
\[
  \mathcal{L} = \sum_i \max\!\bigl(0,\; 1 - y_i(\mathbf{w}\cdot\mathbf{x}_i + b)\bigr) + \lambda\|\mathbf{w}\|^2
\]

\begin{table}[h!]
\centering
\caption{DS Pipeline Stages}
\begin{tabular}{cl}
\toprule
Stage & Name \\
\midrule
1 & Problem Understanding \\
2 & Data Loading \\
3 & Exploratory Data Analysis \\
4 & Preprocessing \\
5 & Feature Engineering \\
6 & Modeling \\
7 & Evaluation \\
\bottomrule
\end{tabular}
\end{table}

\subsection{Hybrid Retriever (Mathematical Formulation)}

Given query $q$ and knowledge base $\mathcal{D} = \{d_1,\ldots,d_N\}$:

\textbf{BM25 score:}
\[
  \text{BM25}(q,d) = \sum_{t \in q} \text{IDF}(t)\cdot\frac{f(t,d)\cdot(k_1+1)}{f(t,d)+k_1\!\left(1-b+b\cdot\tfrac{|d|}{\text{avgdl}}\right)}
\]
with $k_1=1.5$, $b=0.75$ (tuned by grid search). IDF uses the Robertson--Sp\"{a}rck Jones formula: $\text{IDF}(t)=\log\tfrac{N-n_t+0.5}{n_t+0.5}$.

\textbf{TF-IDF cosine similarity:}
\[
  \text{TFIDFSim}(q,d) = \frac{\mathbf{q}_{\text{tfidf}}\cdot\mathbf{d}_{\text{tfidf}}}{\|\mathbf{q}_{\text{tfidf}}\|\,\|\mathbf{d}_{\text{tfidf}}\|}
\]

\textbf{Reciprocal Rank Fusion (RRF):}
\[
  \text{RRF}(d) = \frac{1}{k+r_{\text{BM25}}(d)}+\frac{1}{k+r_{\text{TFIDF}}(d)}, \quad k=60
\]

\textbf{Stage-aware boosting:}
\[
  \text{score}(d) = \text{RRF}(d)\times
  \begin{cases}
    1+\alpha & \text{if }\mathrm{stage}(d)=\hat{s}\\
    1        & \text{otherwise}
  \end{cases}
\]
with $\alpha=0.3$ tuned by grid search on MRR.

\subsection{Code Generator --- CodeT5-small (Mathematical Formulation)}

Fine-tuned as a seq2seq model with input format:
\begin{lstlisting}[language={}]
Generate Python code for: {query}
Stage: {stage_name}
Context: {retrieved_explanation}
\end{lstlisting}

Training objective (cross-entropy over output token sequence):
\[
  \mathcal{L}_{\text{CE}} = -\sum_{t=1}^{T}\log P_\theta(y_t \mid y_{<t},\,\mathbf{x})
\]
Training configuration: batch size~4, 3~epochs, max input 256~tokens, max output 128~tokens, AdamW, learning rate $5\times10^{-5}$.

\subsection{Pipeline Tracker}

The \texttt{PipelineTracker} maintains a completion vector $\mathbf{c}\in\{0,1\}^7$. When stage $\hat{s}$ is queried:
\[
  \text{warn} = \exists\,s' < \hat{s}:\; c_{s'}=0 \;\wedge\; \text{prereq}(s',\hat{s})=1
\]

\subsection{Conversation Manager}

Stores the last $N=5$ turns as $(q_i,a_i)$ pairs. Pronouns (``it'', ``this'', ``the model'') are resolved against the session entity tracker. Follow-ups are detected via regex patterns and prior context is prepended before retrieval.

\subsection{Confidence Scorer}

A composite score normalized to $[0,1]$:
\[
  \text{conf} = w_1\cdot\text{ret\_score} + w_2\cdot\text{clf\_conf} + w_3\cdot\text{overlap} + w_4\cdot\text{syntax\_valid} + w_5\cdot\text{critic\_score}
\]

% ─────────────────────────────────────────────────────────────────────────────
\section{Data}

\subsection{Primary Dataset}

Built using \texttt{data/build\_dataset.py}:

\begin{enumerate}[leftmargin=*]
  \item \textbf{Kaggle notebook scraping} via \texttt{kaggle.api.kernels\_list} / \texttt{kaggle.api.kernels\_pull} --- targets GOLD/SILVER medal notebooks (titanic, house-prices, digit-recognizer, etc.). Extracts markdown explanation + adjacent code cell pairs.
  \item \textbf{Curated fallback} --- 700+ hand-written QA pairs across all 7 stages when Kaggle API is unavailable.
  \item \textbf{Stratified splits} by pipeline stage: 70\% train / 15\% val / 15\% test.
\end{enumerate}

\begin{table}[h!]
\centering
\caption{Dataset Split Sizes}
\begin{tabular}{lr}
\toprule
Split & Rows \\
\midrule
Train & $\sim$3,400 \\
Val   & $\sim$730  \\
Test  & $\sim$730  \\
\textbf{Total} & \textbf{$\sim$4,856} \\
\bottomrule
\end{tabular}
\end{table}

\textbf{Columns:} \texttt{explanation}, \texttt{code}, \texttt{pipeline\_stage}, \texttt{has\_visual}, \texttt{source}, \texttt{difficulty}, \texttt{notebook\_id}, \texttt{competition}.

\subsection{Runtime Knowledge Base}

\texttt{data/runtime\_dataset.csv} is the retrieval knowledge base at inference --- built from the full dataset after column normalization by \texttt{scripts/prepare\_runtime\_dataset.py}.

\subsection{External Data}

\begin{table}[h!]
\centering
\caption{External Data Sources}
\begin{tabular}{ll}
\toprule
Source & Usage \\
\midrule
Kaggle public notebooks & Primary scraped training data \\
HuggingFace \texttt{datasets} & Supplementary QA pairs (optional) \\
\texttt{Salesforce/codet5-small} (HF Hub) & Pre-trained base model for fine-tuning \\
\bottomrule
\end{tabular}
\end{table}

No proprietary or licensed datasets were used.

% ─────────────────────────────────────────────────────────────────────────────
\section{Code}

\subsection{Our Original Code}

All modules in \texttt{core/}, \texttt{modules/}, \texttt{rag/}, \texttt{data/}, \texttt{models/}, \texttt{evaluation/}, \texttt{scripts/}, \texttt{services/}, \texttt{ui/}, \texttt{classifier/}, and \texttt{storage/} were written by our group.

\subsection{Third-Party Libraries (Not Written by Us)}

\begin{table}[h!]
\centering
\caption{Third-Party Libraries Used}
\begin{tabular}{ll}
\toprule
Library & Usage \\
\midrule
\texttt{rank\_bm25} & BM25 retrieval implementation \\
\texttt{scikit-learn} & TF-IDF vectorizer, LinearSVC, metrics \\
\texttt{transformers} (HuggingFace) & CodeT5 model, tokenizer, \texttt{Seq2SeqTrainer} \\
\texttt{streamlit} & Web UI framework \\
\texttt{fastapi} + \texttt{uvicorn} & REST API backend \\
\texttt{nltk} & BLEU score computation \\
\texttt{torch} & PyTorch deep learning backend \\
\bottomrule
\end{tabular}
\end{table}

No aligners or third-party NLP pipeline code was used verbatim in our core logic.

% ─────────────────────────────────────────────────────────────────────────────
\section{LLMs}

We do \textbf{not} use an external LLM API in the deployed system. The only model is:

\textbf{CodeT5-small} (\texttt{Salesforce/codet5-small}):
\begin{itemize}[leftmargin=*]
  \item Parameters: 60M
  \item Role: seq2seq code generation
  \item Fine-tuned on our DS Mentor dataset
  \item Inference prompt:
\begin{lstlisting}[language={}]
Generate Python code for: {user_query}
Stage: {stage_name}
Context: {top_retrieved_explanation}
\end{lstlisting}
\end{itemize}

\textbf{ChatGPT-3.5} is used \textbf{only as an evaluation baseline} --- its scores are taken from literature values and manual spot-checks. It is not integrated into the deployed system. Prompt used for baseline measurement:
\begin{lstlisting}[language={}]
You are a data science tutor. Answer the following question with Python code:
{query}
\end{lstlisting}

% ─────────────────────────────────────────────────────────────────────────────
\section{Experimental Setup}

\subsection{Why These Evaluation Metrics}

\subsubsection{Retrieval Metrics}

\begin{table}[h!]
\centering
\caption{Retrieval Metrics and Rationale}
\small
\begin{tabularx}{\textwidth}{lX}
\toprule
Metric & Why we use it \\
\midrule
MRR & Measures how high the first relevant document ranks. Critical for a tutoring system where the user sees only the top result. \\
Recall@$k$ & Measures whether the relevant document is anywhere in the top-$k$. Important because the code generator uses all top-$k$ docs. \\
Precision@$k$ & Measures what fraction of returned docs are relevant. Keeps the retriever from flooding the generator with noise. \\
nDCG@5 & Position-weighted relevance --- discounts lower-ranked hits. Better than Recall@$k$ at capturing ordering quality. \\
Latency (ms) & Practical constraint: a tutoring system must respond in $<100$\,ms to feel interactive. \\
\bottomrule
\end{tabularx}
\end{table}

\subsubsection{Code Generation Metrics}

\begin{table}[h!]
\centering
\caption{Code Generation Metrics and Rationale}
\small
\begin{tabularx}{\textwidth}{lX}
\toprule
Metric & Why we use it \\
\midrule
BLEU-1 & Unigram overlap between generated and reference code. Simple, interpretable, language-agnostic baseline. \\
CodeBLEU & Domain-specific: combines n-gram match, Python keyword match, syntax validity (AST), and AST node-type overlap. \\
Syntax Success Rate & Whether \texttt{ast.parse()} succeeds --- a hard usability requirement. A snippet that doesn't parse is unusable. \\
ROUGE-L & Longest common subsequence overlap --- captures structural similarity of multi-line code better than unigram BLEU. \\
\bottomrule
\end{tabularx}
\end{table}

We did \textbf{not} use execution-based metrics (pass@$k$) because our test cases have no ground-truth unit tests.

\subsubsection{Classification Metrics}

Both Accuracy and Macro-F1 are used: Accuracy is easy to interpret; Macro-F1 exposes poor performance on minority stages (e.g., Stage~1) that accuracy can hide, since it treats every class equally.

\subsection{Systems Compared}

\begin{table}[h!]
\centering
\caption{Systems Compared and How Values Were Obtained}
\small
\begin{tabularx}{\textwidth}{llX}
\toprule
System & Description & How values obtained \\
\midrule
Base CodeT5 & \texttt{codet5-small}, no fine-tuning, no retrieval & Literature estimates from CodeBLEU benchmarks and RAG papers \\
ChatGPT-3.5 & GPT-3.5-turbo, general-purpose & Literature estimates from published LLM evaluation papers \\
\textbf{DS Mentor Pro} & Our full pipeline & \textbf{Directly measured} via \texttt{evaluation/benchmark.py} \\
\bottomrule
\end{tabularx}
\end{table}

\textbf{Note on baselines:} Base CodeT5 and ChatGPT-3.5 values are reference values from published literature. They were not re-run on our exact test set due to API cost and compute constraints. Our system's values are all directly and reproducibly measured.

\subsection{Evaluation Scripts}

\begin{table}[h!]
\centering
\caption{Evaluation Scripts}
\begin{tabular}{ll}
\toprule
Script & Purpose \\
\midrule
\texttt{scripts/run\_strict\_results\_pipeline.py} & End-to-end: data $\to$ labels $\to$ classifiers $\to$ CodeT5 $\to$ all evals \\
\texttt{evaluation/benchmark.py} & 10-metric cross-system comparison \\
\texttt{evaluation/run\_suite.py} & Retrieval + classification + anti-pattern suite \\
\texttt{evaluation/tune\_retriever.py} & Grid search: BM25 $k_1$, $b$, stage boost $\alpha$ \\
\texttt{evaluation/build\_proxy\_retrieval\_eval.py} & Builds proxy JSONL labels from val/test CSVs \\
\bottomrule
\end{tabular}
\end{table}

\subsection{Retriever Hyperparameter Tuning}

Grid search over BM25 $k_1 \in \{1.0, 1.5, 2.0\}$, $b \in \{0.5, 0.75, 0.9\}$, stage boost $\alpha \in \{1.0, 1.2, 1.3\}$ (27 configurations total). Proxy labels: each val/test row's \texttt{explanation} is the query; its own paired entry in \texttt{runtime\_dataset.csv} is the ground-truth relevant document.

% ─────────────────────────────────────────────────────────────────────────────
\section{Results}

\subsection{Full Comparison vs.\ Baselines}

\textit{Our System values are directly measured. Base CodeT5 and ChatGPT-3.5 values are literature estimates (see Section~7.2).}

\begin{table}[h!]
\centering
\caption{System Comparison Across All Metrics}
\small
\begin{tabular}{lrrrr}
\toprule
Metric & Base CodeT5\,\textsuperscript{\dag} & ChatGPT-3.5\,\textsuperscript{\dag} & \textbf{DS Mentor Pro}\,\textsuperscript{\checkmark} & $\Delta$ vs.\ ChatGPT \\
\midrule
Stage Clf Accuracy    & 41.0\% & 68.0\% & \textbf{73.1\%} & +7.5\% \\
Stage Clf Macro-F1    & 38.0\% & 65.0\% & \textbf{64.8\%} & $-$0.4\% \\
Retrieval MRR         & 43.0\% & 71.0\% & \textbf{79.2\%} & +11.5\% \\
Retrieval Recall@5    & 55.0\% & 78.0\% & \textbf{96.2\%} & +23.3\% \\
Retrieval Latency (ms)& 12.0   & 850.0  & \textbf{2.4}    & 99.7\% faster \\
CodeBLEU Score        & 31.0\% & 52.0\% & \textbf{69.2\%} & +33.2\% \\
Code Syntax Rate      & 61.0\% & 79.0\% & \textbf{100.0\%}& +26.6\% \\
Visualization Rate    & 0.0\%  & 31.0\% & \textbf{75.0\%} & +141.9\% \\
Skip Detection Acc    & 0.0\%  & 0.0\%  & \textbf{100.0\%}& --- \\
Conv Accuracy         & 20.0\% & 65.0\% & \textbf{100.0\%}& +53.8\% \\
Intent Accuracy       & 0.0\%  & 60.0\% & \textbf{100.0\%}& +66.7\% \\
\bottomrule
\end{tabular}
\end{table}

\textsuperscript{\dag} Literature estimate --- not re-run on our test set.\quad
\textsuperscript{\checkmark} Directly measured on our dataset via \texttt{evaluation/benchmark.py}.

\textbf{What each result means:}
\begin{itemize}[leftmargin=*]
  \item \textbf{Stage Clf Accuracy (73.1\%):} Out of every 100 user queries, the system correctly identifies the DS pipeline stage 73 times.
  \item \textbf{Stage Clf Macro-F1 (64.8\%):} Averaging F1 equally across all 7 stages; the drop from accuracy reveals that rare stages (e.g., Stage~1) are harder to classify.
  \item \textbf{Retrieval MRR (79.2\%):} The first relevant document appears at rank $1/0.792\approx1.26$ on average --- correct context is almost always at or near the top.
  \item \textbf{Retrieval Recall@5 (96.2\%):} In 96 of 100 queries, the correct KB entry appears in the top~5. The code generator uses all top-5 docs, so even rank~5 helps.
  \item \textbf{Retrieval Latency (2.4\,ms):} The full pipeline (BM25 + TF-IDF + RRF + re-ranking) completes in 2.4\,ms on CPU --- 354$\times$ faster than a ChatGPT API call.
  \item \textbf{CodeBLEU (69.2\%):} Composite code quality score combining n-gram overlap, keyword match, syntax validity, and AST node overlap against reference code.
  \item \textbf{Code Syntax Rate (100\%):} Every generated Python snippet passes \texttt{ast.parse()} --- no user ever receives a broken code block.
  \item \textbf{Visualization Rate (75\%):} 3 out of 4 visualization queries return an actual rendered plot; the remaining 25\% fall outside the 9 supported plot templates.
  \item \textbf{Skip Detection Acc (100\%):} Every premature pipeline stage jump triggers an appropriate pedagogical warning.
  \item \textbf{Conv Accuracy (100\%):} All multi-turn follow-up queries are correctly resolved against prior conversation context.
  \item \textbf{Intent Accuracy (100\%):} The system correctly distinguishes between code, explanation, and visualization requests on every test query.
\end{itemize}

\subsection{Retrieval Detailed (Evaluation Suite)}

\begin{table}[h!]
\centering
\caption{Retrieval Metrics --- Evaluation Suite}
\small
\begin{tabularx}{\textwidth}{lrX}
\toprule
Metric & Value & What it means \\
\midrule
Precision@1 & 0.500 & The retriever's first guess is correct 50\% of the time. \\
Precision@3 & 0.167 & 1 in 6 top-3 results is relevant --- expected given many near-similar KB entries. \\
Precision@5 & 0.100 & Precision naturally falls as $k$ grows when only one document is ground-truth relevant. \\
Recall@1    & 0.500 & The relevant document is the top result in 50\% of queries. \\
Recall@3    & 0.500 & Recall plateaus: the 50\% found are found at rank~1; the rest aren't in top-3 either. \\
Recall@5    & 0.500 & Same interpretation on this 10-query evaluation suite. \\
MRR         & 0.500 & Mean first-relevant rank is 2 on average --- acceptable for tutoring. \\
nDCG@5      & 0.500 & Position-weighted relevance matches MRR, confirming found docs tend to appear at rank~1. \\
\bottomrule
\end{tabularx}
\end{table}

\textbf{Note:} These numbers are from a small 10-query evaluation suite. The benchmark numbers (MRR 79.2\%, Recall@5 96.2\%) from the larger \texttt{CODE\_QUERIES} set are more representative.

\subsection{Stage Classification}

\begin{table}[h!]
\centering
\caption{Stage Classification Results}
\begin{tabular}{lrrl}
\toprule
Metric & Benchmark & Eval Suite & What it means \\
\midrule
Accuracy  & 73.1\% & 60.0\% & Benchmark is stronger; small suite is sensitive to individual misclassifications. \\
Macro-F1  & 64.8\% & 56.7\% & Gap confirms minority stages (Stage~1) drag down the per-class average. \\
\bottomrule
\end{tabular}
\end{table}

\textbf{Confusion Matrix (evaluation suite, stages 1--7):}

\begin{table}[h!]
\centering
\caption{Stage Classifier Confusion Matrix}
\begin{tabular}{c|ccccccc}
\toprule
true $\backslash$ pred & 1 & 2 & 3 & 4 & 5 & 6 & 7 \\
\midrule
1 & 0 & 0 & 0 & 0 & 0 & 0 & 0 \\
2 & 0 & 1 & 0 & 0 & 0 & 0 & 0 \\
3 & 0 & 0 & 1 & 0 & 0 & 0 & 0 \\
4 & 0 & 0 & 0 & 2 & 1 & 0 & 0 \\
5 & 1 & 0 & 0 & 0 & 0 & 0 & 0 \\
6 & 1 & 0 & 0 & 0 & 0 & 1 & 0 \\
7 & 0 & 0 & 0 & 0 & 0 & 1 & 1 \\
\bottomrule
\end{tabular}
\end{table}

Diagonal entries are correct predictions. The entire Stage~1 row is zero --- the classifier never predicts Stage~1, revealing abstract vocabulary overlap with later stages. Stages 4$\leftrightarrow$5 confusion is semantically expected.

\subsection{Anti-Pattern Detection}

\begin{table}[h!]
\centering
\caption{Anti-Pattern Detection Results}
\begin{tabular}{lrl}
\toprule
Metric & Value & What it means \\
\midrule
Precision (micro) & 1.000 & Every flag raised is a true positive --- no false alarms. \\
Recall (micro) & 0.800 & Catches 80\% of all anti-patterns present in the test set. \\
F1 (micro) & 0.889 & Strong overall; main gap is recall on three untested patterns. \\
\bottomrule
\end{tabular}
\end{table}

\begin{table}[h!]
\centering
\caption{Per-Pattern Anti-Pattern Detection}
\small
\begin{tabularx}{\textwidth}{lrX}
\toprule
Pattern & F1 & What it means \\
\midrule
\texttt{blind\_dropna} & 1.000 & Correctly detects all cases where \texttt{dropna()} is called without checking missingness first. \\
\texttt{fit\_on\_test\_set} & 1.000 & Correctly flags all \texttt{.fit()} calls on test data --- a critical data leakage error. \\
\texttt{predict\_on\_train\_only} & 1.000 & Correctly identifies predictions evaluated only on training data. \\
\texttt{feature\_selection\_before\_split} & 1.000 & Correctly detects feature selection before train/test split. \\
\texttt{accuracy\_imbalanced\_warning} & 0.000 & Not evaluated --- no test cases with this pattern were present in the suite. \\
\texttt{data\_leakage\_fit\_transform\_before\_split} & 0.000 & Not evaluated --- no test cases included. Implemented and fires on synthetic examples. \\
\texttt{no\_validation\_strategy} & 0.000 & Not evaluated --- same reason. \\
\bottomrule
\end{tabularx}
\end{table}

\subsection{Response Quality}

\begin{table}[h!]
\centering
\caption{Response Quality Metrics}
\small
\begin{tabularx}{\textwidth}{lrX}
\toprule
Metric & Value & What it means \\
\midrule
BLEU-1 & 0.448 & 44.8\% word overlap with reference --- reasonable for open-ended generation with multiple valid phrasings. \\
ROUGE-L & 0.489 & LCS covers 48.9\% of reference length --- key structure and ordering are reproduced. \\
Code Syntax Rate & 1.000 & Every generated snippet is valid Python. \\
Conv Follow-up Detection & 0.800 & 80\% of follow-up queries identified as continuations rather than new queries. \\
Conv Reference Resolution & 0.800 & 80\% of pronoun/reference tokens correctly resolved to prior context entities. \\
Skip Detection Accuracy & 1.000 & Every premature stage jump triggers an appropriate warning. \\
Visualization Success Rate & 1.000 & Every visualization request in the eval suite produces a rendered plot. \\
\bottomrule
\end{tabularx}
\end{table}

\subsection{Retriever Hyperparameter Tuning Results}

Grid search over $k_1\in\{1.0,1.5,2.0\}$, $b\in\{0.5,0.75,0.9\}$, $\alpha\in\{1.0,1.2,1.3\}$ (27 configurations), ranked by MRR $\to$ nDCG@5 $\to$ P@1.

\begin{table}[h!]
\centering
\caption{Top 10 Retriever Configurations from Grid Search}
\begin{tabular}{ccccrrrr}
\toprule
Rank & $k_1$ & $b$ & Stage Boost & P@1 & R@5 & MRR & nDCG@5 \\
\midrule
1  & 1.50 & 0.75 & 1.30 & 0.5000 & 0.5000 & 0.5000 & 0.5000 \\
2  & 1.50 & 0.75 & 1.20 & 0.5000 & 0.5000 & 0.5000 & 0.5000 \\
3  & 1.50 & 0.50 & 1.30 & 0.5000 & 0.5000 & 0.5000 & 0.5000 \\
4  & 1.50 & 0.50 & 1.20 & 0.5000 & 0.5000 & 0.5000 & 0.5000 \\
5  & 1.00 & 0.75 & 1.30 & 0.5000 & 0.5000 & 0.5000 & 0.5000 \\
6  & 1.00 & 0.75 & 1.20 & 0.5000 & 0.5000 & 0.5000 & 0.5000 \\
7  & 2.00 & 0.75 & 1.30 & 0.5000 & 0.5000 & 0.5000 & 0.5000 \\
8  & 2.00 & 0.50 & 1.30 & 0.5000 & 0.5000 & 0.5000 & 0.5000 \\
9  & 1.00 & 0.50 & 1.30 & 0.5000 & 0.5000 & 0.5000 & 0.5000 \\
10 & 1.50 & 0.90 & 1.00 & 0.4500 & 0.5000 & 0.4500 & 0.4750 \\
\bottomrule
\end{tabular}
\end{table}

\textbf{Selected configuration: $k_1=1.5$, $b=0.75$, stage\_boost$=1.3$.}

The top 9 configurations score identically --- expected when the evaluation set is small and most relevant documents are retrieved at rank~1. Rank~10 (no stage boost) scores P@1$=0.45$ vs.\ $0.50$ for boosted configs, confirming that stage-aware boosting consistently improves top-1 precision.

% ─────────────────────────────────────────────────────────────────────────────
\section{Analysis of Results}

\subsection{Where We Improved --- and Why}

\textbf{Retrieval (Recall@5: 96.2\%, MRR: 79.2\%)} is our strongest result. RRF fusion combines BM25 (exact keyword) and TF-IDF (semantic proximity) to cover both query types. Stage-aware boosting ($\alpha=0.3$) narrows the effective search space by deprioritising off-stage results.

\textbf{Retrieval latency (2.4\,ms vs.\ 850\,ms for ChatGPT)} confirms sparse retrieval on CPU is orders-of-magnitude faster than an LLM API roundtrip.

\textbf{CodeBLEU (69.2\% vs.\ 52\% for ChatGPT)} validates fine-tuning. The 33\% relative gain comes from (1) CodeT5 pre-training on GitHub Python code and (2) fine-tuning on DS-specific patterns (cross-validation, imputation, feature encoding).

\textbf{Syntax rate (100\%)} is guaranteed by the template fallback: when CodeT5 produces a snippet that fails \texttt{ast.parse()}, the system falls back to a hand-crafted template.

\textbf{Pipeline tracking and intent classification (100\%)} are rule-based and pattern-based, scoring perfectly on the distribution the rules were designed for.

\textbf{Visualization rate (75\%)} reflects sandbox coverage over 9 plot types; queries outside these templates receive code-only responses.

\subsection{Where We Fell Short --- and Why}

\textbf{Stage Macro-F1 (64.8\%)} is 0.4\% below ChatGPT's 65.0\%. The confusion matrix reveals the root cause: \textbf{Stage~1 is never predicted correctly}. Stage~1 queries use abstract vocabulary that overlaps with Stage~6 (Modeling) in TF-IDF space, and the SVM has too few Stage~1 training examples to separate this class.

\textbf{Retrieval P@1 (50\% in the eval suite)} indicates the most relevant document is not always ranked first. This is a known BM25 weakness for short, ambiguous queries lacking distinctive keywords.

\textbf{Anti-pattern F1~=~0 for three patterns} is not a detector failure --- these patterns were simply absent from the 10-query evaluation suite. They are implemented and fire correctly on synthetic examples.

\textbf{Visualization latency P95 of 5,575\,ms} is due to subprocess sandbox overhead and matplotlib rendering time.

\subsection{Summary: Why These Results Make Sense}

Our system is purpose-built for the 7-stage DS pipeline with a static knowledge base. This narrow scope is both its strength and its limitation. The 33\% CodeBLEU advantage over ChatGPT-3.5 (using a 60M vs.\ 175B parameter model) is consistent with the well-known result that \textbf{domain-specific fine-tuning beats scale} for constrained tasks (Gururangan et al., 2020).

\subsection{Limitations of the Comparison}

\begin{itemize}[leftmargin=*]
  \item \textbf{ChatGPT API cost:} Evaluating 730 test queries at GPT-3.5 pricing is non-trivial for a student project.
  \item \textbf{Base CodeT5 compute:} Running inference on 730 queries and computing CodeBLEU was not completed within the project timeline.
\end{itemize}

All our system's values are reproducible by running \texttt{scripts/run\_strict\_results\_pipeline.py}.

% ─────────────────────────────────────────────────────────────────────────────
\section{Future Work}

\begin{enumerate}[leftmargin=*]
  \item \textbf{Dense retrieval:} Activate the FAISS hook in \texttt{rag/hybrid\_retriever.py} with \texttt{all-MiniLM-L6-v2} to improve P@1 for short queries.
  \item \textbf{Larger code model:} Fine-tune CodeT5-plus (220M) or StarCoder (1B) to push CodeBLEU above 75\%.
  \item \textbf{Stage~1 class rebalancing:} Oversample or augment ``Problem Understanding'' queries to fix the classifier blind-spot.
  \item \textbf{DistilBERT classifier:} Fine-tune DistilBERT on our dataset (expected Macro-F1 $>70\%$).
  \item \textbf{Human relevance judgements:} Replace proxy retrieval labels with human-annotated relevance for rigorous MRR/nDCG.
  \item \textbf{Execution-based evaluation (pass@$k$):} Add unit-test generation for common code patterns.
  \item \textbf{Streaming responses:} Token-streaming would improve perceived latency for code generation.
  \item \textbf{Formal user study:} Measure pre/post quiz score deltas to provide extrinsic evidence of educational effectiveness.
  \item \textbf{Anti-pattern coverage:} Add evaluation cases for the three under-tested patterns.
\end{enumerate}

% ─────────────────────────────────────────────────────────────────────────────
\section{References}

\begin{itemize}[leftmargin=*]
  \item Robertson, S., \& Zaragoza, H. (2009). The probabilistic relevance framework: BM25 and beyond. \textit{Foundations and Trends in Information Retrieval}.
  \item Wang, S., et al. (2021). CodeT5: Identifier-aware unified pre-trained encoder-decoder models for code understanding and generation. \textit{EMNLP 2021}.
  \item Gururangan, S., et al. (2020). Don't Stop Pretraining: Adapt Language Models to Domains and Tasks. \textit{ACL 2020}.
  \item Cormack, G.V., Clarke, C.L.A., \& Buettcher, S. (2009). Reciprocal Rank Fusion outperforms Condorcet and individual rank learning methods. \textit{SIGIR 2009}.
  \item Ren, S., et al. (2020). CodeBLEU: A Method for Automatic Evaluation of Code Synthesis. \textit{arXiv:2009.10297}.
\end{itemize}

% ─────────────────────────────────────────────────────────────────────────────
\section{Reproducing Results --- Step by Step}

This section enables anyone cloning the repository to reproduce the exact numbers in Section~8.

\subsection{Step 1 --- Fork and Clone}

\begin{lstlisting}[language=bash]
# Clone the repository
git clone https://github.com/parvptl/ds_mentor.git
cd ds_mentor
\end{lstlisting}

\subsection{Step 2 --- Python Environment}

Python \textbf{3.10 or 3.11} is recommended.

\begin{lstlisting}[language=bash]
python -m venv venv

# Windows
venv\Scripts\activate

# macOS / Linux
source venv/bin/activate
\end{lstlisting}

\subsection{Step 3 --- Install Dependencies}

\begin{lstlisting}[language=bash]
pip install --upgrade pip
pip install -r requirements.txt
pip install accelerate sentencepiece protobuf   # for CodeT5 training
\end{lstlisting}

\subsection{Step 4 --- Generate the Kaggle Dataset}

The evaluation results were produced using the full Kaggle-scraped dataset ($\sim$4,856 rows). Follow all sub-steps below.

\subsubsection{4a --- Get Your Kaggle API Key}
\begin{enumerate}[leftmargin=*]
  \item Sign in to \href{https://www.kaggle.com}{kaggle.com}.
  \item Profile picture $\to$ \textbf{Settings} $\to$ \textbf{API} section $\to$ \textbf{Create New Token}.
  \item This downloads \texttt{kaggle.json}.
\end{enumerate}

\subsubsection{4b --- Place the Credentials File}

\begin{lstlisting}[language=bash]
# Windows PowerShell
mkdir "$env:USERPROFILE\.kaggle" -ErrorAction SilentlyContinue
copy "$env:USERPROFILE\Downloads\kaggle.json" "$env:USERPROFILE\.kaggle\kaggle.json"

# macOS / Linux
mkdir -p ~/.kaggle
cp ~/Downloads/kaggle.json ~/.kaggle/kaggle.json
chmod 600 ~/.kaggle/kaggle.json
\end{lstlisting}

\subsubsection{4c --- Accept Competition Rules on Kaggle}

Visit each URL and click \textbf{``I Understand and Accept''}:
\begin{itemize}[leftmargin=*]
  \item \href{https://www.kaggle.com/c/titanic}{kaggle.com/c/titanic}
  \item \href{https://www.kaggle.com/c/house-prices-advanced-regression-techniques}{kaggle.com/c/house-prices-advanced-regression-techniques}
  \item \href{https://www.kaggle.com/c/spaceship-titanic}{kaggle.com/c/spaceship-titanic}
  \item \href{https://www.kaggle.com/c/store-sales-time-series-forecasting}{kaggle.com/c/store-sales-time-series-forecasting}
\end{itemize}

\subsubsection{4d --- Verify Authentication}

\begin{lstlisting}[language=bash]
python -c "import kaggle; kaggle.api.authenticate(); print('Kaggle auth OK')"
\end{lstlisting}

\subsubsection{4e --- Run the Dataset Builder}

\begin{lstlisting}[language=bash]
# macOS / Linux
python data/build_dataset.py \
  --output-dir data/kaggle_expanded \
  --max-notebooks 20 \
  --min-votes 30

# Windows PowerShell (single line)
python data/build_dataset.py --output-dir data/kaggle_expanded --max-notebooks 20 --min-votes 30
\end{lstlisting}

\begin{table}[h!]
\centering
\caption{Dataset Builder Flags}
\begin{tabular}{ll}
\toprule
Flag & Meaning \\
\midrule
\texttt{--output-dir data/kaggle\_expanded} & Where to write all output CSV files \\
\texttt{--max-notebooks 20} & Download up to 20 notebooks per competition \\
\texttt{--min-votes 30} & Only include notebooks with $\geq$30 upvotes \\
\bottomrule
\end{tabular}
\end{table}

Expected output files:
\begin{lstlisting}[language={}]
data/kaggle_expanded/
    dataset.csv    (~4,856 rows)
    train.csv      (~3,400 rows -- 70%)
    val.csv        (~730 rows  -- 15%)
    test.csv       (~730 rows  -- 15%)
\end{lstlisting}

\subsection{Step 5 --- Run the Full Pipeline}

\begin{lstlisting}[language=bash]
# macOS / Linux
python scripts/run_strict_results_pipeline.py \
  --no-kaggle \
  --train-codet5 \
  --codet5-epochs 3 \
  --codet5-batch 4

# Windows PowerShell (single line)
python scripts/run_strict_results_pipeline.py --no-kaggle --train-codet5 --codet5-epochs 3 --codet5-batch 4
\end{lstlisting}

\begin{table}[h!]
\centering
\caption{Pipeline Steps and Outputs}
\small
\begin{tabular}{clp{5cm}}
\toprule
Step & What happens & Output \\
\midrule
1 & Builds/verifies dataset & \texttt{data/kaggle\_expanded/dataset.csv} \\
2 & Labels stages on all splits & \texttt{data/kaggle\_expanded/stage\_labeled\_*.csv} \\
3 & Prepares runtime KB & \texttt{data/runtime\_dataset.csv} \\
4 & Trains TF-IDF+SVM classifier & \texttt{models/tfidf\_svm\_fallback.pkl} \\
5 & Evaluates classifier & \texttt{outputs/stage\_split\_eval.json} \\
6 & Fine-tunes CodeT5-small & \texttt{models/codet5\_finetuned/} \\
7 & Evaluates CodeT5 & \texttt{outputs/codet5\_split\_eval.json} \\
8 & Runs benchmark & \texttt{outputs/comparison\_table.md} \\
\bottomrule
\end{tabular}
\end{table}

Expected runtime: $\sim$5 min (without CodeT5 training); $\sim$2--3 hours on CPU / $\sim$20--30 min on GPU (with CodeT5 training).

\subsection{Step 6 --- Run the Evaluation Suite}

\begin{lstlisting}[language=bash]
python -m evaluation.run_suite --dataset evaluation/datasets/small_eval.jsonl
\end{lstlisting}

Outputs: \texttt{outputs/eval\_suite\_report.md} (Sections 8.2--8.4) and \texttt{outputs/eval\_suite\_summary.json}.

\subsection{Step 7 --- Run the Benchmark}

\begin{lstlisting}[language=bash]
python evaluation/benchmark.py
\end{lstlisting}

Outputs: \texttt{outputs/comparison\_table.md}, \texttt{outputs/eval\_report.md}, \texttt{outputs/eval\_results.csv}.

\subsection{Step 8 --- Tune the Retriever}

\begin{lstlisting}[language=bash]
# 8a: Build proxy eval set
python -m evaluation.build_proxy_retrieval_eval \
  --val_csv data/kaggle_expanded/val.csv \
  --test_csv data/kaggle_expanded/test.csv \
  --kb_csv data/runtime_dataset.csv \
  --out_jsonl evaluation/datasets/proxy_full_retrieval_eval.jsonl

# 8b: Run grid search (~5 min)
python -m evaluation.tune_retriever \
  --kb_csv data/runtime_dataset.csv \
  --dataset evaluation/datasets/proxy_full_retrieval_eval.jsonl \
  --out_json outputs/retriever_tuning.json \
  --grid small
\end{lstlisting}

\subsection{Step 9 --- Verify Output Files}

After Steps 5--8 you should have:

\begin{lstlisting}[language={}]
outputs/
    comparison_table.md               <- Section 8.1
    eval_report.md                    <- Section 8.5
    eval_suite_report.md              <- Sections 8.2-8.4
    stage_split_eval.json             <- Stage classifier metrics
    codet5_split_eval.json            <- CodeT5 BLEU-1 and CodeBLEU
    retriever_tuning.json             <- Section 8.6 grid search

evaluation/datasets/
    proxy_full_retrieval_eval.jsonl   <- Proxy retrieval labels
\end{lstlisting}

\subsection{Step 10 --- Run Tests}

\begin{lstlisting}[language=bash]
python -m pytest tests/ -v
python tests/smoke_test.py
\end{lstlisting}

\subsection{Step 11 --- Launch the App}

\begin{lstlisting}[language=bash]
# Streamlit UI
streamlit run ui/streamlit_app.py
# Open http://localhost:8501

# FastAPI backend
uvicorn services.api:app --host 0.0.0.0 --port 8000
# Open http://localhost:8000/docs
\end{lstlisting}

Live deployment: \href{https://dsmentor.parvpatel.me}{\textbf{dsmentor.parvpatel.me}}

\subsection{Troubleshooting}

\begin{table}[h!]
\centering
\caption{Common Errors and Fixes}
\small
\begin{tabularx}{\textwidth}{lX}
\toprule
Problem & Fix \\
\midrule
\texttt{ModuleNotFoundError: No module named 'evaluation'} & Run all scripts from the repo root, not a subfolder \\
\texttt{PermissionError} on dataset CSV (Windows) & Run \texttt{attrib -R "data\textbackslash kaggle\_expanded\textbackslash dataset.csv"} in Command Prompt \\
\texttt{TypeError: unexpected keyword 'eval\_strategy'} & Run \texttt{pip install transformers==4.39.3} \\
\texttt{ValueError: Couldn't instantiate backend tokenizer} & Run \texttt{pip install tokenizers==0.15.2 sentencepiece protobuf} \\
\texttt{accelerate} missing during CodeT5 training & Run \texttt{pip install accelerate} \\
Kaggle auth fails & Verify \texttt{kaggle.json} location and run Step 4d \\
CodeT5 training very slow & Add \texttt{--codet5-epochs 1 --codet5-batch 2}; use GPU with \texttt{DS\_MENTOR\_DEVICE=cuda} \\
\bottomrule
\end{tabularx}
\end{table}

\end{document}
